% StreamingLLM: Efficient Streaming Language Models with Attention Sinks
% Consolidated from text/ directory

\section{Abstract}
Deploying Large Language Models (LLMs) in streaming applications such as multi-round dialogue poses two major challenges:
\begin{enumerate}
\item During decoding, caching previous tokens' Key and Value states (KV) consumes extensive memory
\item Popular LLMs cannot generalize to longer texts than the training sequence length
\end{enumerate}

Window attention (caching only most recent KVs) fails when text length surpasses cache size. We observe an interesting phenomenon called "attention sink" - keeping KV of initial tokens largely recovers window attention performance.

StreamingLLM enables LLMs trained with finite length attention window to generalize to infinite sequence length without fine-tuning. It enables Llama-2, MPT, Falcon, and Pythia to perform stable language modeling with up to 4 million tokens. StreamingLLM achieves up to 22.2x speedup over sliding window recomputation baseline.

\section{Introduction}

\subsection{Problem}
LLMs are constrained by the attention window during pre-training (e.g., 4K for Llama-2). Two primary challenges for infinite input streams:
\begin{enumerate}
\item KV cache leads to excessive memory usage and increasing decoding latency
\item Limited length extrapolation - performance degrades when sequence length exceeds pre-training window
\end{enumerate}

\subsection{Window Attention Failure}
Window attention maintains fixed-size sliding window on KV states but collapses once sequence length exceeds cache size - even just evicting the KV of the first token causes failure.

\subsection{Attention Sink Discovery}
Key observation: A surprisingly large amount of attention score is allocated to initial tokens, irrespective of their relevance to the language modeling task. We term these "attention sinks".

Despite lacking semantic significance, initial tokens collect significant attention scores because:
\begin{itemize}
\item SoftMax requires attention scores to sum to one
\item Model dumps unnecessary attention values to specific tokens
\item Initial tokens are visible to almost all subsequent tokens (autoregressive nature)
\item Makes them readily trained to serve as attention sinks
\end{itemize}

\subsection{StreamingLLM Solution}
Keep attention sink tokens' KV (just 4 initial tokens suffice) together with sliding window's KV to anchor attention computation and stabilize performance.

\section{Method}

\subsection{Failure Analysis of Window Attention}
\textbf{Perplexity Surge:} Perplexity spikes when text length surpasses cache size due to exclusion of initial tokens.

\textbf{Why initial tokens matter:} Beyond bottom two layers, model consistently focuses on initial tokens across all layers and heads. Removing initial tokens' KV removes considerable portion of SoftMax denominator:
\begin{equation}
\text{SoftMax}(x)_i = \frac{e^{x_i}}{e^{x_1} + \sum_{j=2}^N e^{x_j}}, \quad x_1 \gg x_j
\end{equation}

\textbf{Position vs Semantics:} Experiments replacing first four tokens with linebreak "\textbackslash n" show model still emphasizes these tokens. The absolute position of starting tokens, rather than semantic value, holds greater significance.

\subsection{Rolling KV Cache with Attention Sinks}
KV cache divided into two parts:
\begin{enumerate}
\item Attention sinks (four initial tokens) - stabilize attention computation
\item Rolling KV Cache - retains most recent tokens for language modeling
\end{enumerate}

\textbf{Position Encoding:} Focus on positions within the cache rather than in original text. For cache [0,1,2,3,6,7,8] decoding 9th token, positions assigned are [0,1,2,3,4,5,6,7] not [0,1,2,3,6,7,8,9].

For RoPE: Cache Keys prior to rotary transformation, apply position transformation at each decoding phase.
For ALiBi: Apply contiguous linear bias instead of 'jumping' bias.

\subsection{Pre-Training with Sink Tokens}
Include global trainable attention sink token at start of all training samples. Alternative: SoftMax-off-by-One:
\begin{equation}
\text{SoftMax}_1(x)_i = \frac{e^{x_i}}{1 + \sum_{j=1}^N e^{x_j}}
\end{equation}
(Equivalent to prepending token with all-zero Key and Value features)

\textbf{Results:} Model trained with sink token achieves satisfactory streaming performance using just the sink token, while vanilla model requires multiple initial tokens.

\section{Experiments}

\subsection{Models Evaluated}
Llama-2, MPT, Pythia, Falcon families with RoPE and ALiBi position encodings.

\subsection{Language Modeling Results}
\begin{itemize}
\item StreamingLLM matches oracle baseline (sliding window with re-computation) on 20K token texts
\item Reliably handles 4+ million tokens across model families and scales
\item Dense attention fails when input exceeds pre-training window
\item Window attention fails when input exceeds cache size
\end{itemize}

\subsection{Pre-Training with Sink Token Results}
\begin{itemize}
\item No negative impact on model convergence or NLP benchmark performance
\item Similar performance on ARC, HellaSwag, LAMBADA, OpenbookQA, PIQA, Winogrande
\item Vanilla model needs multiple initial tokens; sink token model needs only the sink
\end{itemize}

\subsection{Streaming Question Answering}
Using Llama-2 Chat models on ARC datasets:
\begin{itemize}
\item Dense attention: Out-of-Memory errors
\item Window attention: Low accuracy due to random outputs
\item StreamingLLM: Matches one-shot baseline accuracy
\end{itemize}

StreamEval benchmark (query every 10 lines, answer 20 lines prior):
\begin{itemize}
\item StreamingLLM maintains accuracy even approaching 120K tokens
\item Can complement context extension techniques
\end{itemize}

\subsection{Ablation Studies}
\textbf{Number of Initial Tokens:}
\begin{itemize}
\item 1-2 initial tokens insufficient
\item 4 initial tokens threshold is enough
\item Additional tokens have marginal effects
\end{itemize}

\textbf{Cache Sizes:} Increasing cache size doesn't consistently lower perplexity - models may not maximize utility of entire context.

\subsection{Efficiency Results}
\begin{itemize}
\item StreamingLLM: Linear growth in decoding speed with cache size
\item Re-computation baseline: Quadratic rise in latency
\item Up to 22.2x speedup per token
\item Consistent memory footprint with baseline
\end{itemize}

\section{Key Takeaways}
\begin{enumerate}
\item Attention sinks are initial tokens that collect unnecessary attention due to SoftMax constraints
\item StreamingLLM keeps 4 initial tokens + sliding window for stable infinite-length generation
\item Pre-training with dedicated sink token eliminates need for multiple initial tokens
\item 22.2x speedup over re-computation baseline with same quality
\item Works with RoPE and ALiBi position encodings
\end{enumerate}
